% !TEX root = ../main.tex
% !TEX program = lualatex
\documentclass[../main.tex]{subfiles}
\begin{document}
\IfSubfilesClassLoaded{\chapter{EDO Study Guide}}{}

\section{Availability Work Package Requirements and Development Stages}
The \ac{jfmm} defines the \ac{awp} as the consolidated industrial work plan for a scheduled maintenance availability, combining maintenance, modernization, and repair requirements drawn from the \ac{cmp}/\ac{icmp}, authorized alterations, \ac{tycom} repairs, pre-availability tests, and the \ac{csmp}.~\autocite{COMFLTFORCOMINST4790-3,edo-4-2-1-cno-availability-planning-2025} The \ac{awp} progresses through sequential stages (baseline, preliminary, proposed, approved, and completed), with baseline planning starting about three years before the availability and surface-ship baseline packages developed by \ac{surfmepp}.~\autocite{COMFLTFORCOMINST4790-3,edo-4-2-1-cno-availability-planning-2025} Submarine baseline planning is managed by \ac{submepp}, and the completed \ac{awp} documents work accomplished and lessons learned within months after availability completion.~\autocite{COMFLTFORCOMINST4790-3,edo-4-2-1-cno-availability-planning-2025}

\section{Planning Milestones and Availability Type Differences}
Planned milestones for \ac{sra} and other \ac{cno}-scheduled availabilities establish early schedule alignment, work package integration, and contracting actions, with key gates at roughly A-720 (availability schedule), A-360 (letter of authorization), A-220 (work package integration conference), A-165 (100 percent work package lock letter), A-120 (contract award), and A-45 to A-30 (work package execution review) before A-0 start of availability.~\autocite{COMFLTFORCOMINST4790-3,edo-4-2-1-cno-availability-planning-2025,edo-4-2-2-cno-availability-execution-2025} By contrast, \ac{cmav} planning compresses the cycle, with planning beginning around A-180, integration and ship checks closer to A-60 to A-20, and a shortened execution review near A-10 before A-0 start.~\autocite{COMFLTFORCOMINST4790-3,edo-4-2-1-cno-availability-planning-2025,edo-4-2-2-cno-availability-execution-2025} Both paths use progress conferences during execution and culminate in completion reporting and lessons learned capture after availability closeout.~\autocite{edo-4-2-2-cno-availability-execution-2025}

\section{System Operational Verification Test Definition and Significance}
The \ac{sovt} is the structured verification that an alteration or ship change performs as intended following maintenance, and \ac{ait} or \ac{isea} teams must conduct these tests under conditions that simulate normal service.~\autocite{TS9090-310} The alteration is not complete until the \ac{sovt} or stage testing is successfully executed and witnessed by ship's force, and test reports must be provided to ship's force, the \ac{nsa}, and the \ac{lma}.~\autocite{TS9090-310} \ac{sovt} completion is distinct from the \ac{pcd}, so schedules must protect time after production completion to execute and document testing.~\autocite{TS9090-310}

\section{Signed and Unsigned System Operational Verification Tests}
A signed \ac{sovt} is a completed test plan that has been executed, witnessed, and recorded with results and discrepancies, enabling formal acceptance and closeout of the alteration.~\autocite{TS9090-310} An unsigned \ac{sovt} indicates the test was not fully executed or witnessed, or that discrepancies remain unresolved, which delays acceptance and introduces schedule risk through retesting or deferred closeout actions.~\autocite{TS9090-310}

Current as of 2026-02-02.

\end{document}
